\subsection{Conditional ranges under the original hump}\label{app:shapeprofile}
Using all 256 standard and same-expiry midcurve ATM observations gives feasible sets conditional on the original hump $(\chi,\kappa)=(4,0.75)$. The median minimum ATM-price allowance drops from 3.319 bp under parallel loading to 1.136 bp with the hump and constant background, 1.078 bp with 90-day cells, and 0.936 bp with expiry-gap cells. At a 1 bp allowance, these three humped specifications fit respectively 2, 5, and 9 of the fourteen dates; no parallel specification fits a date.

On 22 July, the first date fitting all three humped backgrounds at that allowance, the September-meeting ranges are $[88,550]$, $[0,550]$, and $[0,1038]$ bp$^2$ for constant, 90-day, and gap backgrounds. The January 2027 ranges are $[414,973]$, $[378,2113]$, and $[314,2514]$ bp$^2$. Same-expiry windows now add background information: with expiry-gap cells the matrix rank rises from 7 under parallel loading to 12 out of 15 columns under the hump. In contrast to the standard-only parallel experiment, assigning each increment to its own background cell no longer makes all meeting lower bounds vanish.

Conditioning on an estimated shape can overstate this information. Taking the union of the constant-background feasible sets over the entire grid gives September range $[0,550]$ and January range $[0.16,973]$ bp$^2$ on 22 July. Only two grid shapes fit that date within 1 bp. The tiny positive January endpoint is specific to the finite grid; it is not a positive bound over all smooth loadings. These original-hump calculations illustrate the effect of conditioning on a shape. They are separate from the free-long-end prediction comparison in Section~\ref{sec:humpfit}.

